\begin{proof}[Proof of \cref{obj:thm:generic-apolar-arrow-recovery}]
Put \(K=2m+2\) and \(n=m+2\).  For a direction \(u=(u_1,u_2)\in\mathbb C^2\) write
\[
\ell_u(x,y)=u_1x+u_2y,
\qquad
q(u)=q(u_1,u_2)\ \text{the evaluation of a binary form }q\text{ at }u,
\]
and, for a loading family \((w_j)_{j=0}^{m+1}\), write its support annihilator in the normalization of \cref{obj:def:apolar-notation} as the degree-\(n\) form
\[
Q_D=\prod_{j=0}^{m+1}\bigl(w_{j2}\,x-w_{j1}\,y\bigr),
\]
one linear factor per loading direction, each vanishing exactly on that direction.  Throughout, \(\Theta^{\mathrm{right}}_{m,K}\) and \(\Theta^{\mathrm{left}}_{m,K}\) are the finite retained-coordinate parameter spaces of \cref{obj:def:truncated-cumulant}, in which every source weight outside the orders \(2\le r\le K\) is zero.

\begin{enumerate}
\item \textbf{Contraction of a divided-power block.}  For a forward parameter \(\theta=(\gamma,\rho,c)\) and \(2\le r\le K\), expanding the binomial in \(\ell_{u_j}^r\) and comparing with \cref{obj:def:forward-cumulant-map} gives
\[
f_r(x,y)=\sum_{a=0}^{r}\binom{r}{a}\bigl[\Phi^{\mathrm{right}}_{m,K}(\theta)\bigr]_{r,a}x^{r-a}y^{a}
=\sum_{j=0}^{m+1}c_{jr}\,\ell_{u_j}(x,y)^{r},
\]
and the same identity with \(d_{jr}\), \(v_j\), and \(\Phi^{\mathrm{left}}_{m,K}\) follows from \cref{obj:def:reverse-cumulant-map}.  If \(q\) is a homogeneous binary form of degree \(n\), then differentiating a power of a linear form gives, for every \(k\ge0\),
\[
q(\partial)\,\ell_u^{\,n+k}=\frac{(n+k)!}{k!}\;q(u)\;\ell_u^{\,k}.
\]
Consequently, for \(0\le k\le m\),
\[
q(\partial)f_{n+k}
=\frac{(n+k)!}{k!}\sum_{j=0}^{m+1}c_{j,n+k}\,q(u_j)\,\ell_{u_j}^{\,k},
\]
and the scalar \((n+k)!/k!\) is nonzero.
% lean: dividedPowerBlock dividedPowerBlock_forward_eq_sum_linForm_pow diffApply diffApply_linForm_pow

\item \textbf{The contraction minor and the loci.}  For a forward parameter define the \((m+2)\times(m+2)\) matrix \(M^{\mathrm{right}}(\theta)\) with rows indexed by \(i=0,\ldots,m+1\) and columns by the source labels \(j=0,\ldots,m+1\),
\[
M^{\mathrm{right}}(\theta)_{0,j}=c_{j,n},
\qquad
M^{\mathrm{right}}(\theta)_{a+1,j}=\binom{m}{a}c_{j,K}\,(u_{j1})^{m-a}(u_{j2})^{a}
\qquad(0\le a\le m),
\]
and let \(M_{\mathrm r}(\theta)=\det M^{\mathrm{right}}(\theta)\), a polynomial in the coordinates of \(\theta\).  These rows are the degree-zero block and the coefficients of the degree-\(m\) block in the weighted contraction map
\[
e=(e_j)_{j=0}^{m+1}\ \longmapsto
\Bigl(\sum_{j=0}^{m+1}c_{j,n+k}\,e_j\,\ell_{u_j}^{\,k}\Bigr)_{k=0}^{m}.
\]
Thus nonsingularity of this selected square coefficient matrix forces the displayed weighted contraction map to be injective.

For a reverse parameter define instead
\[
M^{\mathrm{left}}(\eta)_{0,j}=d_{j,n},
\qquad
M^{\mathrm{left}}(\eta)_{a+1,j}=\binom{m}{a}d_{j,K}\,(v_{j1})^{a}(v_{j2})^{m-a}
\qquad(0\le a\le m),
\]
and let \(M_{\mathrm l}(\eta)=\det M^{\mathrm{left}}(\eta)\).  The last \(m+1\) rows are the coefficients of \(x^ay^{m-a}\) in the degree-\(m\) reverse contraction block, so \(M_{\mathrm l}(\eta)\ne0\) gives injectivity of the corresponding reverse weighted contraction map.

Each of \(M_{\mathrm r}\) and \(M_{\mathrm l}\) is a nonzero polynomial.  At the forward parameter
\[
\theta^\star:\quad \gamma=1,
\quad \rho_i=i+1\ (1\le i\le m),
\quad c_{j,n}=\mathbb 1_{\{j=m+1\}},
\quad c_{j,K}=\mathbb 1_{\{j\le m\}},
\quad c_{jr}=0\ \text{otherwise},
\]
the matrix \(M^{\mathrm{right}}(\theta^\star)\) has last column \((1,0,\ldots,0)^{\top}\) and, in its remaining columns and last \(m+1\) rows, the entries \(\binom{m}{a}(j+1)^{a}\); it is therefore triangular with a Vandermonde determinant in the distinct nodes \(1,\ldots,m+1\) as its surviving factor, so \(M_{\mathrm r}(\theta^\star)\ne0\).  All weights of \(\theta^\star\) vanish outside the orders \(2\le r\le K\), so \(M_{\mathrm r}\) remains a nonzero polynomial after every off-band weight coordinate is set to \(0\).  The reverse witness has \(\delta=m+1\), \(\sigma_i=i+1\), \(d_{0,n}=1\), \(d_{j,K}=1\) for \(j\ne0\), and all other retained witness weights zero; the same Vandermonde calculation with the displayed reverse minor gives the corresponding statements for \(M_{\mathrm l}\).
% lean: forwardWeightedContraction contraction_injective_of_minor_det_ne_zero forward_contraction_injective_of_generic_and_minor reverse_contraction_injective_of_generic_and_minor witnessParameter

\item \textbf{The cut-out polynomials.}  Let \(\Delta_{\mathrm{gen}}^{\mathrm r}\) and \(\Delta_{\mathrm{gen}}^{\mathrm l}\) be the forward and reverse generic polynomials from \cref{obj:def:generic-parameter-loci}, namely
\[
\Delta_{\mathrm{gen}}^{\mathrm r}(\gamma,\rho,c)
=\gamma\prod_{i=1}^m(\gamma-\rho_i)\prod_{1\le i<i'\le m}(\rho_i-\rho_{i'})
\prod_{j=0}^{m+1}\prod_{r=2}^{K}c_{jr},
\]
and
\[
\Delta_{\mathrm{gen}}^{\mathrm l}(\delta,\sigma,d)
=\delta\prod_{i=1}^m(\delta-\sigma_i)\prod_{1\le i<i'\le m}(\sigma_i-\sigma_{i'})
\prod_{j=0}^{m+1}\prod_{r=2}^{K}d_{jr}.
\]
Put \(\Delta_\rho=\prod_{i=1}^m\rho_i\) in the forward chart and \(\Delta_\sigma=\prod_{i=1}^m\sigma_i\) in the reverse chart, and define
\[
R_{\mathrm r}=\Delta_{\mathrm{gen}}^{\mathrm r}\,\Delta_\rho\,M_{\mathrm r},
\qquad
R_{\mathrm l}=\Delta_{\mathrm{gen}}^{\mathrm l}\,\Delta_\sigma\,M_{\mathrm l}.
\]
Each factor is a nonzero polynomial that stays nonzero after the off-band weight coordinates are set to \(0\): this is immediate for \(\Delta_{\mathrm{gen}}^{\mathrm r}\), \(\Delta_{\mathrm{gen}}^{\mathrm l}\), \(\Delta_\rho\), and \(\Delta_\sigma\), since they involve only slopes and retained weights, and it follows for the determinants from the witness calculations above.  Hence \(R_{\mathrm r}\) and \(R_{\mathrm l}\) have the same two properties.
% lean: genericParameterPolynomial genericParameterLocus_eq_nonvanishing_poly pinSubst_ne_zero_of_pinned_witness

\item \textbf{The loci \(U^{\mathrm{right}}_m\) and \(U^{\mathrm{left}}_m\).}  Set
\[
U^{\mathrm{right}}_m=\bigl\{\theta\in\Theta^{\mathrm{right}}_{m,K}:R_{\mathrm r}(\theta)\ne0\bigr\},
\qquad
U^{\mathrm{left}}_m=\bigl\{\eta\in\Theta^{\mathrm{left}}_{m,K}:R_{\mathrm l}(\eta)\ne0\bigr\}.
\]
Inside the retained-coordinate space, the nonvanishing locus of a polynomial that is not identically zero there is the complement of a proper relatively closed set, hence relatively Zariski open, and it is relatively Zariski dense because any polynomial vanishing on it multiplies \(R_{\mathrm r}\), respectively \(R_{\mathrm l}\), to the zero polynomial and therefore vanishes identically.  Since \(\Delta_{\mathrm{gen}}^{\mathrm r}\) divides \(R_{\mathrm r}\), we get \(U^{\mathrm{right}}_m\subseteq\Theta^{\mathrm{right},\circ}_{m,K}\); likewise \(\Delta_{\mathrm{gen}}^{\mathrm l}\) divides \(R_{\mathrm l}\), so \(U^{\mathrm{left}}_m\subseteq\Theta^{\mathrm{left},\circ}_{m,K}\).
% lean: isZariskiOpenIn_denseIn_of_pinSubst_ne_zero IsZariskiOpenParamIn IsZariskiDenseParamIn

\item \textbf{Open neighborhoods meeting feasibility.}  Set
\[
O^{\mathrm{right}}=\{\theta\in\mathbb R^{m+1}\times\mathbb R^{(m+2)(K-1)}:R_{\mathrm r}(\theta_{\mathbb C})\ne0\},
\qquad
O^{\mathrm{left}}=\{\eta\in\mathbb R^{m+1}\times\mathbb R^{(m+2)(K-1)}:R_{\mathrm l}(\eta_{\mathbb C})\ne0\},
\]
with \(\theta_{\mathbb C}\) and \(\eta_{\mathbb C}\) the coordinatewise complexifications of \cref{obj:def:apolar-notation}.  These are Euclidean open because \(\theta\mapsto R_{\mathrm r}(\theta_{\mathbb C})\) and \(\eta\mapsto R_{\mathrm l}(\eta_{\mathbb C})\) are continuous.  The truncated moment problem supplies a real cumulant vector \(c^{0}\) and a radius \(\varepsilon>0\) such that
\[
|c'_r-c^{0}_r|<\varepsilon\ \ (2\le r\le K)
\quad\Longrightarrow\quad
\exists\ \text{a centered non-Gaussian law }\nu\text{ with } \mathbb E|\,\cdot\,|^{K}<\infty
\text{ and }\operatorname{Cum}_r(\nu)=c'_r\ (2\le r\le K).
\]
Consider the box in which each slope coordinate ranges over all of \(\mathbb R\) and each retained weight coordinate \(c_{jr}\) ranges over \((c^{0}_r-\varepsilon,c^{0}_r+\varepsilon)\).  Each factor of the box is infinite, and \(R_{\mathrm r}\) is a nonzero polynomial after pinning the off-band coordinates, so it cannot vanish at every point of the box; pick a point \(\theta\) where it does not vanish.  Then \(\Delta_{\mathrm{gen}}^{\mathrm r}(\theta)\ne0\) gives \(\gamma\ne0\) and pairwise distinct slopes, and each weight family \(c_{j\cdot}\) lies in the \(\varepsilon\)-box, hence is realized by a centered non-Gaussian real source law with finite \(K\)-th moment.  Thus \(\theta\in F^{\mathrm{right}}_{m,K}\) by \cref{obj:def:real-feasible-regions}, and \(R_{\mathrm r}(\theta_{\mathbb C})\ne0\) puts \(\theta\in O^{\mathrm{right}}\).  Hence \(O^{\mathrm{right}}\cap F^{\mathrm{right}}_{m,K}\ne\varnothing\), and the same argument with \(R_{\mathrm l}\) gives \(O^{\mathrm{left}}\cap F^{\mathrm{left}}_{m,K}\ne\varnothing\).  Finally, if \(\theta\in O^{\mathrm{right}}\cap F^{\mathrm{right}}_{m,K}\), then \(\theta_{\mathbb C}\) has all off-band weights zero because feasibility pins them, and \(R_{\mathrm r}(\theta_{\mathbb C})\ne0\), so \(\theta_{\mathbb C}\in U^{\mathrm{right}}_m\); likewise on the reverse side:
\[
\bigl(O^{\mathrm{right}}\cap F^{\mathrm{right}}_{m,K}\bigr)_{\mathbb C}\subseteq U^{\mathrm{right}}_m,
\qquad
\bigl(O^{\mathrm{left}}\cap F^{\mathrm{left}}_{m,K}\bigr)_{\mathbb C}\subseteq U^{\mathrm{left}}_m .
\]
% lean: TruncatedMomentInterior truncatedMomentInterior exists_feasible_nonvanishing isOpen_realNonvanishingLocus

\item \textbf{What a point of \(U^{\mathrm{right}}_m\) provides.}  Let \(\theta=(\gamma,\rho,c)\in U^{\mathrm{right}}_m\).  Factoring \(R_{\mathrm r}(\theta)\ne0\) gives
\[
\Delta_{\mathrm{gen}}^{\mathrm r}(\theta)\ne0,
\qquad \rho_i\ne0\ (1\le i\le m),
\qquad M_{\mathrm r}(\theta)\ne0 .
\]
Hence \(\theta\) is generic, \(\gamma\ne0\), the second coordinates of the \(m+1\) finite loading directions,
\[
u_{02}=\gamma,
\qquad
u_{j2}=\rho_j\ (1\le j\le m),
\]
are pairwise distinct and all nonzero, and the forward weighted contraction map displayed above is injective.
% lean: forward_data gamma_ne_zero_of_generic forward_slopes_injective_of_generic forward_slopes_ne_zero_of_generic

\item \textbf{The common contraction kernel is the line \(\mathbb C\,Q_D\).}  Let \(D\) be the forward loading-direction support of \(\theta\) and
\[
Q_D=\bigl(\gamma x-y\bigr)\prod_{j=1}^{m}\bigl(\rho_jx-y\bigr)\cdot x,
\]
the last factor coming from \(u_{m+1}=(0,1)\).  If \(q=cQ_D\) is homogeneous of degree \(n\), then \(q(u_j)=0\) for every \(j\), so every contraction \(q(\partial)f_{n+k}\), \(0\le k\le m\), vanishes by the divided-power contraction identity.  Conversely, suppose \(q\) is homogeneous of degree \(n\) and \(q(\partial)f_{n+k}=0\) for all \(0\le k\le m\).  Dividing the contraction identity by the nonzero factorial, the vector \(e=(q(u_j))_{j=0}^{m+1}\) is sent to \(0\) by the forward weighted contraction map.  Injectivity therefore gives
\[
q(u_j)=0\qquad\text{for every }j=0,\ldots,m+1 .
\]
The \(n=m+2\) directions \(u_0,\ldots,u_{m+1}\) are pairwise non-proportional, because their second coordinates \(\gamma,\rho_1,\ldots,\rho_m\) are distinct and nonzero and \(u_{m+1}\) is the vertical direction.  A degree-\(n\) binary form vanishing at these \(n\) pairwise non-proportional directions is divisible by the corresponding product of linear forms, i.e. by \(Q_D\); comparing degrees gives \(q=cQ_D\) for some \(c\in\mathbb C\).  Thus, for every homogeneous binary form \(q\) of degree \(n\),
\[
\bigl[q(\partial)f_{n+k}=0\ \text{for every }0\le k\le m\bigr]
\quad\Longleftrightarrow\quad
\exists c\in\mathbb C:\ q=cQ_D .
\]
Moreover \(Q_D\ne0\), every factor being a nonzero linear form; \(x\mid Q_D\) by the displayed factorization; and \(y\nmid Q_D\), since the factor \(x\) is not divisible by \(y\), and the factors \(\gamma x-y\), \(\rho_jx-y\) are not because \(\gamma\ne0\) and \(\rho_j\ne0\).
% lean: supportAnnihilator forward_supportAnnihilator_in_contraction_kernel evalAtDir_zero_of_forward_contractions forward_points_imply_supportAnnihilator_multiple forward_apolar_kernel_identity supportAnnihilator_ne_zero X0_dvd_supportAnnihilator_forward not_X1_dvd_supportAnnihilator_forward

\item \textbf{The reverse fiber over \(\Phi^{\mathrm{right}}_{m,K}(\theta)\) is empty.}  Suppose \(\eta=(\delta,\sigma,d)\in R^{\mathrm{left}}_{m,K}\!\bigl(\Phi^{\mathrm{right}}_{m,K}(\theta)\bigr)\).  Fiber membership in \cref{obj:def:fiber-correspondences} compares the two maps on the retained coordinates, and both cumulant maps vanish off that range, so
\[
\Phi^{\mathrm{left}}_{m,K}(\eta)=\Phi^{\mathrm{right}}_{m,K}(\theta),
\qquad\text{hence}\qquad
f^{\mathrm{left}}_r=f_r\ \text{ for all }r .
\]
Let \(Q_D^{\mathrm{left}}=\prod_j(v_{j2}x-v_{j1}y)\) be the support annihilator of the reverse loading family of \(\eta\).  It is homogeneous of degree \(n\), it is nonzero, and it vanishes at every reverse direction.  Therefore it annihilates all reverse blocks \(f^{\mathrm{left}}_{n+k}\), \(0\le k\le m\), by the divided-power contraction identity, and hence annihilates all forward blocks \(f_{n+k}\).  The forward kernel identity then yields \(c\in\mathbb C\) with
\[
Q_D^{\mathrm{left}}=c\,Q_D .
\]
If \(c=0\), this makes \(Q_D^{\mathrm{left}}=0\), a contradiction.  If \(c\ne0\), then \(y\mid Q_D^{\mathrm{left}}\), because the reverse family contains \(v_0=(1,0)\) and therefore contributes the factor \(-y\); the displayed scalar multiple then forces \(y\mid Q_D\), contradicting the forward divisibility conclusion above.  Hence
\[
R^{\mathrm{left}}_{m,K}\!\left(\Phi^{\mathrm{right}}_{m,K}(\theta)\right)=\varnothing .
\]
% lean: forward_reverse_fiber_empty reverse_supportAnnihilator_in_contraction_kernel X1_dvd_supportAnnihilator_reverse

\item \textbf{Recovery of the slope multiset.}  Put
\[
Q(z)=(z-\gamma)\prod_{i=1}^{m}(z-\rho_i).
\]
The entries \(\gamma,\rho_1,\ldots,\rho_m\) are pairwise distinct and nonzero, so \(Q\ne0\), \(Q\) is squarefree, \(Q(0)=(-1)^{m+1}\gamma\prod_i\rho_i\ne0\), and the multiset of roots of \(Q\) is exactly \(\{\gamma,\rho_1,\ldots,\rho_m\}\).  Up to sign, \(Q\) is the dehomogenization \(Q_D(1,z)\): the factors \(\gamma x-y\) and \(\rho_jx-y\) become \(\gamma-z\) and \(\rho_j-z\), and the factor \(x\) becomes \(1\).  Now let \(\theta''\) be a generic forward parameter with
\[
\Phi^{\mathrm{right}}_{m,K}(\theta'')=\Phi^{\mathrm{right}}_{m,K}(\theta),
\]
and let \(Q_D''\) be its support annihilator.  Its own support annihilator annihilates its own divided-power blocks, hence the blocks of \(\theta\); the forward kernel identity gives \(Q_D''=c\,Q_D\) with \(c\ne0\), because \(Q_D''\ne0\).  Dehomogenizing at \(x=1\) gives \(Q_D''(1,z)=c\,Q_D(1,z)\), so the two dehomogenizations have the same root multiset.  Therefore
\[
\{\gamma'',\rho''_1,\ldots,\rho''_m\}=\{\gamma,\rho_1,\ldots,\rho_m\}.
\]
% lean: qDefault qDefault_roots qDefault_squarefree qDefault_eval_zero_ne rtsF_nodup zero_notMem_rtsF forward_slopes_determined_by_cumulants roots_dehomX_supportAnnihilator_forward

\item \textbf{The reverse locus.}  The argument is the same after exchanging the two coordinate axes, hence the roles of \(x\) and \(y\) and of the two arrows.  For \(\eta=(\delta,\sigma,d)\in U^{\mathrm{left}}_m\), factoring \(R_{\mathrm l}(\eta)\ne0\) gives that \(\eta\) is generic, \(\delta\ne0\), all \(\sigma_i\ne0\), the first coordinates
\[
v_{j1}=\sigma_j\ (1\le j\le m),
\qquad
v_{m+1,1}=\delta
\]
of the finite reverse loading directions are pairwise distinct and nonzero, and the reverse weighted contraction map is injective.  Writing \(Q_D^{\mathrm{left}}\) for the support annihilator of the reverse loading family of \(\eta\),
\[
Q_D^{\mathrm{left}}=y\cdot\prod_{j=1}^{m}\bigl(x-\sigma_jy\bigr)\cdot\bigl(x-\delta y\bigr)
\]
up to sign, the first factor coming from \(v_0=(1,0)\).  Hence \(Q_D^{\mathrm{left}}\ne0\), \(y\mid Q_D^{\mathrm{left}}\), and \(x\nmid Q_D^{\mathrm{left}}\), because \(\delta\ne0\) and all \(\sigma_i\ne0\).  Repeating the divided-power contraction and interpolation argument with the reverse map gives the reverse kernel identity, for \(t=\Phi^{\mathrm{left}}_{m,K}(\eta)\) and \(f_r\) built from \(t\):
\[
\bigl[q(\partial)f_{n+k}=0\ \text{for every }0\le k\le m\bigr]
\quad\Longleftrightarrow\quad
\exists c\in\mathbb C:\ q=c\,Q_D^{\mathrm{left}} .
\]
A forward parameter \(\theta\in R^{\mathrm{right}}_{m,K}\bigl(\Phi^{\mathrm{left}}_{m,K}(\eta)\bigr)\) would have its own support annihilator \(Q_D^{\mathrm{right}}\) in that kernel, hence \(Q_D^{\mathrm{right}}=c\,Q_D^{\mathrm{left}}\) with \(c\ne0\).  But \(x\mid Q_D^{\mathrm{right}}\), because the forward family contains \(u_{m+1}=(0,1)\), while \(x\nmid Q_D^{\mathrm{left}}\), a contradiction.  Therefore
\[
R^{\mathrm{right}}_{m,K}\!\left(\Phi^{\mathrm{left}}_{m,K}(\eta)\right)=\varnothing .
\]
Finally \(Q(z)=(z-\delta)\prod_{i=1}^m(z-\sigma_i)\) is nonzero, squarefree, satisfies \(Q(0)\ne0\), and has root multiset \(\{\delta,\sigma_1,\ldots,\sigma_m\}\).  Dehomogenizing the reverse kernel identity at \(y=1\) shows that every generic reverse parameter with the same cumulants has that same multiset.
% lean: reverse_data reverse_apolar_kernel_identity reverse_forward_fiber_empty X1_dvd_supportAnnihilator_reverse not_X0_dvd_supportAnnihilator_reverse X0_dvd_supportAnnihilator_forward reverse_slopes_determined_by_cumulants roots_dehomY_supportAnnihilator_reverse
\end{enumerate}
\end{proof}
